\section{Results}
\label{sec:results}

\subsection{Placement has a crossover, not a winner}
\label{sec:placement-results}

The controlled family first isolates decision dimension.  At $n=64$,
exact GPU-AD/CPU-KKT takes 22.6\,ms versus 43.1\,ms for DLPack GPU KKT.  At
$n=784$, the ordering reverses: DLPack takes 95.2\,ms, host-staged GPU KKT
136.1\,ms, and CPU KKT 198.8\,ms.  DLPack wraps a 32-MiB JAX buffer in about
40\,$\mu$s, compared with 10.9\,ms for a staged transfer.  Removing transfer
is measurable, but at small $n$ it cannot amortize solver initialization.

The established PFR family produces the same crossover under real constraint
structure (Table~\ref{tab:pfr-scale}).  Case~1 favors CPU KKT; case~2 is near
the boundary; and case~3 favors DLPack/cuDSS by
\PFRCaseThreeSpeedup$\times$.  The case-3 result is stable over five warm
runs: \PFRCaseThreeGPU\,s median (IQR
\PFRCaseThreeGPULo--\PFRCaseThreeGPUHi\,s) versus
\PFRCaseThreeCPU\,s (IQR
\PFRCaseThreeCPULo--\PFRCaseThreeCPUHi\,s).  Same-oracle Ipopt takes
\PFRCaseThreeIpopt\,s median.

\begin{table}[t]
\centering
\caption{Published PFR first-step warm times.  All numeric entries pass the
common residual audit and reproduce the released first control.  A dash is an
audited failure, not a missing fast run.}
\label{tab:pfr-scale}
\small
\setlength{\tabcolsep}{3.5pt}
\begin{tabular}{lrrrr}
\toprule
Case & Exact CPU & Exact DLPack & Ipopt exact & L-BFGS CPU\\
\midrule
1 ($m=400$)   & \textbf{0.043} & 0.069 & 0.044 & 0.385\\
2 ($m=800$)   & 0.204 & \textbf{0.130} & 0.499 & --\\
3 ($m=7{,}500$) & 19.71 & \textbf{2.27} & 27.84 & 530.36\\
\bottomrule
\end{tabular}
\end{table}

\begin{figure}[t]
  \centering
  \includegraphics[width=\linewidth]{pfr_iclr.pdf}
  \caption{Published PFR-NMPC results.  \textbf{Left:} declaring the true
  derivative structure reverses the case-1 hardware ranking.  \textbf{Right:}
  preserving that structure while increasing the published horizon/state
  regime yields a CPU-to-GPU KKT crossover.}
  \label{fig:pfr}
\end{figure}

\subsection{Disclosing structure can dominate moving hardware}
\label{sec:sparsity-results}

The initial PFR port declared dense derivative coordinates.  On case~1 this
changes CPU exact from \PFRCaseOneSparseCPU\,ms to
\PFRCaseOneDenseCPU\,s---a 13.6$\times$ slowdown---and makes GPU KKT appear
preferable.  With the true 5,200/3,049 Jacobian/Hessian entries, CPU exact and
same-oracle Ipopt are essentially tied, and DLPack GPU is slower.  Thus the
hardware conclusion reverses without changing a single function value.

OMLT full space provides a formulation baseline.  At horizon 40 it expands
case~1 to 15,133 variables and 15,100 constraints and takes 21.2\,s, versus
43.3\,ms for the structure-aware gray box.  Its horizon-5 solution agrees in
objective to $4.6\times10^{-11}$.  With maintained OMLT 1.1/1.2 in the
matched software environment, reduced-space construction does not finish
within 180\,s even at horizon 5, so it has no timed entry.

Sequence behavior adds another qualification.  In the 100-step case-1 closed
loop, CPU exact obtains 100/100 common-KKT success with 29.3\,ms median,
30.4\,ms mean, and 77.5\,ms maximum.  A persistent GPU solver has a lower
24.8\,ms median, but one reuse failure triggers a rebuild, raising its mean to
62.1\,ms and maximum to 3.24\,s.  Both control trajectories stay within
$3.6\times10^{-5}$ state infinity error of the published mechanistic
trajectory.  Median-only reporting would select the less predictable
deployment and hide this tail.

The fixed replay separates initialization from persistence
(Appendix Table~\ref{tab:warm-start}).  All \PFRSequenceSuccesses/
\PFRSequenceTotal{} solves pass.  Cold CPU/GPU medians are
\PFRWarmCPUColdIter/\PFRWarmGPUColdIter{} iterations and shifted-primal medians
are \PFRWarmCPUPrimalIter/\PFRWarmGPUPrimalIter; shifted equality multipliers
change no count.  Setting $\mu_0=10^{-7}$ instead gives
\PFRWarmCPUBarrierIter/\PFRWarmGPUBarrierIter{} iterations and
\PFRWarmCPUBarrierTime/\PFRWarmGPUBarrierTime{} ms.  The setpoint change still
needs 9/10 iterations, and persistent GPU reuse still rebuilds once.  Thus
barrier initialization, not equality-dual reuse, removes the median iteration
gap; CPU KKT remains faster for this small sparse system.

\subsection{Curvature should be routed by failure and factorization cost}
\label{sec:curvature-results}

On twelve prospective public-Darcy instances and five starts at each of two
latent resolutions, compact L-BFGS reaches common KKT on
\DarcyRouteBase/\DarcyRouteTotal{} solves.  The five failures occur on three
instances.  Exact DLPack solves all 30 starts from those difficult instances.
Retrying only the five failed starts therefore reaches
\DarcyRouteTotal/\DarcyRouteTotal{} success while adding
\DarcyRouteOverhead\% to total warm time relative to L-BFGS alone.  The policy
was fixed on separate pilot instances, so this is a prospective routing result
rather than a retrospective oracle.

PFR case~3 shows why per-iteration cost is insufficient.  Compact L-BFGS avoids every
Hessian callback but needs 271 iterations and \PFRCaseThreeLBFGS\,s; 516.2\,s
is spent in linear solves.  Exact DLPack needs 11 iterations and
\PFRCaseThreeGPU\,s.  When each iteration still factorizes a 7,500-constraint
KKT system, the extra factorizations dominate the saved Hessian callbacks.

Gauss--Newton occupies a third regime.  On replayable Darcy it is faster than
exact curvature and uses fewer iterations, but explicit residual-Jacobian
formation peaks at 33.5\,GiB---the same capacity class as the exact Hessian.
It is a time approximation, not a memory-free one.  Appendix~\ref{app:curvature}
also shows that reducing L-BFGS history from 20 to 6 increases controlled
iterations from 21--35 to 64--218; history 50 yields only marginal gains over
20.  ``L-BFGS'' is itself not a parameter-free baseline.

\subsection{A surrogate KKT point is not a simulator solution}
\label{sec:validity-results}

The public Darcy experiment already signals optimization shift.  At latent 16,
the median coefficient error is 0.320 and the best representable error is
0.214; at latent 32 these values are 0.311 and 0.149.  More capacity
slightly improves recovery while nearly doubling the gap to the
parameterization floor.  Meanwhile, the optimized-to-true-coefficient FNO
residual ratio falls from 0.85 to 0.42.  Surrogate residual improves much more
than coefficient recovery.

\begin{figure}[t]
  \centering
  \includegraphics[width=\linewidth]{darcy_iclr.pdf}
  \caption{Optimization-aware Darcy validation.  \textbf{Left:} public-data
  coefficient error versus the best latent representation.  \textbf{Right:}
  on the independently replayable prospective split, lower FNO residual does
  not reliably imply lower sparse-PDE residual.  Markers are successful
  instance--start solves and large symbols are descriptive method medians;
  paired inference uses physical instances as the sampling unit.}
  \label{fig:darcy}
\end{figure}

Independent replay turns that signal into a direct prospective test.  Across
\ReplaySolves\ successful method/start solves on twelve disjoint instances,
\ReplayOutfitsTruth\ optimized FNO residuals are below the FNO residual at the true coefficient.
At latent 32,
Gauss--Newton is fastest (\ReplayGNTime\,s) and has the lowest median FNO
residual (\ReplayGNSurrogate\%), but its PDE residual is
\ReplayGNPDE\%.  L-BFGS is slower (\ReplayLBFGSTime\,s) and has a worse FNO
residual (\ReplayLBFGSSurrogate\%), yet its PDE residual is only
\ReplayLBFGSPDE\%.  Exact DLPack gives
\ReplayExactSurrogate\%/\ReplayExactPDE\% FNO/PDE residual.  The method
ranking agrees under surrogate and simulator in only
\ReplayAgreement/\ReplayGroups{} start-level groups.  After aggregating starts
within each physical instance, Gauss--Newton reduces FNO
residual relative to L-BFGS on all twelve instances by a median
\ReplayPairedFNOReduction\ percentage points (95\% bootstrap CI
[\ReplayPairedFNOReductionLow, \ReplayPairedFNOReductionHigh], Wilcoxon
$p=\ReplayPairedFNOPValue$).  The paired PDE change is instead
$+\ReplayPairedPDEDelta$ points with CI
[\ReplayPairedPDEDeltaLow, \ReplayPairedPDEDeltaHigh] and
$p=\ReplayPairedPDEPValue$.  Thus the evidence is a failure of surrogate gains
to transfer consistently, not proof that L-BFGS is physically superior.

An exploratory tenfold increase in Tikhonov regularization does not improve the
mismatch in this ablation: at latent 32 it raises median PDE residual from 2.37\% to 7.67\% for
L-BFGS.  Staying close to an uninformative prior is not a substitute for model
management or high-fidelity validation.

\subsection{Nonconvex timing curves require local-solution context}
\label{sec:mnist-results}

On MNIST, all methods satisfy the same raw network, slack, target, and output
bound audit, but reach different KKT points.  At 5.01M parameters, the
published start gives objective 4.644 for exact CPU, about 4.736 for exact
DLPack/L-BFGS/Ipopt-exact, and 5.625 for Ipopt limited memory.  Over five shared
starts, all \MNISTSolved/\MNISTTotal{} MadNLP solves pass; exact CPU and exact
DLPack have median objective \MNISTExactMedianObjective, while L-BFGS has
\MNISTLBFGSMedianObjective.  L-BFGS remains fastest (0.47\,s median versus
1.31\,s DLPack and 2.98\,s CPU exact), but all four perturbed starts enter
worse basins.  A speed plot without objective values would reverse the useful
interpretation.

\subsection{Memory and cold start define deployable regimes}

Exact latent-32 Darcy peaks at \DarcyExactMemory\,GiB GPU memory versus
\DarcyLBFGSMemory\,GiB for L-BFGS.  The final dense
$1024\times1024$ Hessian is only 8\,MiB: differentiated FNO intermediates, not
the returned matrix, set the capacity threshold.  PFR case~3 exact DLPack
peaks at 43.9\,GiB versus 1.14\,GiB for CPU KKT.  Conversely, its GPU
factorization supplies the decisive speedup.  Capacity and time therefore need
separate axes.

Parameter count is not a sufficient capacity proxy either.  With FNO size
matched within \FNOShapeMaxParameterMismatch\%, changing depth from one to six
layers moves the $32^2$ Hessian peak from \FNOShapeMinMemory{} to
\FNOShapeMaxMemory{} GiB and warm evaluation from \FNOShapeFastHessian{} to
\FNOShapeSlowHessian{} ms (Appendix Table~\ref{tab:topology-memory}).  Random
and trained four-layer graphs both use \FNOShapeTrainedMemory{} GiB, isolating
topology rather than learned values.

Fresh-process execution requires approximately 60--75\,s for the first JAX/cuDSS
specialization.  The case-3 DLPack cold solve takes 71.7\,s despite a
\PFRCaseThreeGPU\,s warm median.  The GPU recommendation applies to persistent
workers or repeated solves, not one-shot service calls.
